Now, we would like to combine the results of the different tasks and see if they give a coherent picture of the logistical map. 
For all three different tasks there is data on $r_1=3.1$ available and thus this case will be used for comparison.

For the symmetrical case ($r_1 = r_2 = r$) we expect a fixed point at $x = y = 1 - \frac{1}{r}$, which in this case corresponds to 
$x = y = 1 - \frac{1}{3.1} = 0.67741935 $. Looking at the period-1 points in Figure \ref{fig:p_orbits_1} we can see that the $r_1 = r_2 = 3.1$ period-1 fixed point is exactly at the 
expected $x = y = 0.67741935$. This matches our theoretical fixed point very well. However, we do not see this fixed point in the corresponding phase portrait (Figure \ref{fig:phaseport-3.1-3.1}) or 
the bifurcation cascade (Figure \ref{fig:bifurcationcascade}). The reason for this is that the fixed point does not seem to be stable or attractive. No starting values of x and y converge to this fixed point. 
Instead, both the phase portrait and the bifurcation cascade show a convergence towards two different points ($x = y = 0.76456652$ and $x=y=0.55801413$). For the bifurcation 
cascade this is represented in the two strings, while for the phase portrait we can see two yellow data points. As we can see in Figure \ref{fig:p_orbits_2}, these two points are the 
period-2 points for these system parameters. The period-2 points are stable and attractive. Higher order orbits were not observed for these parameters.

Next, we can look at the asymmetrical case where $r_1 = 3.1 \ne r_2$. There are several interesting values of $r_2$ that show different behaviour. At $r_2 =3.4$ there are 
still two branches in the bifurcation cascade, which hints at stable and attractive period-2 points. In this case, however, the points are not symmetrical. As we can tell from 
Figure \ref{fig:p_orbits_2} the coordinates of the fixed points are [0.79737624 0.82726612] and [0.49635601 0.49035224]. As $r_2$ is now larger than $r_1$, the y coordinates of the period-2 points moved 
faster away from the period-2 points for $r_1 = r_2 = 3.1 $ than the x coordinates. These stable period-2 points and their asymmetric nature can also be seen in the phase portrait in Figure 
\ref{fig:phaseport-3.1}.

At $r_2=3.8$ the bifurcation cascade already exhibits four different branches. This is in agreement with the period-4 point that exists at this $r_2$ value (Figure \ref{fig:p_orbits_4}). Again, this seems 
to be the stable and attractive fixed point instead of the period-2 or period-1 fixed points.

At $r_2 = 3.95$ the bifurcation cascade shows a point cloud. Similarly, the phase portrait (Figure \ref{fig:phaseport-3.1-3.95}) exhibits multiple different and well distributed points at higher step numbers. 
Interestingly, there seems to be a linear relation between the x and y value. Looking at the phase portrait for $r_1 = 3.85, r_2=3.95$ this relation is not given at higher $r_1$ 
values. Instead, for each x value there is a cloud of y values.

All in all, the four tasks show a coherent picture. In the symmetrical case there is a fixed point at $x = y = 1 - \frac{1}{3.1} = 0.67741935 $, which is, however, unstable and thus 
not visible in the bifurcation cascade or the phase portrait. Instead, the logistical map converges towards period-2 points. At higher $r_2$ we exhibit stable and attractive 
period-4 points and, at even higher values, a point cloud. In all cases, the results of the p-Orbits, the bifurcation cascade and the phase portraits match.   